\section{Reconstruction dimensionnelle et conservation extérieure
}
\label{nuclear:11}
\subsection{Deux mémoires récupèrent le registre centré
}

\subsubsection{Opérateurs de lecture du registre
}

Dans le système de coordonnées ordonné du registre \(K\) de
\cref{sec:k-registro,img09:imagen}, soit \(V=\RR^{12}\) sa
réalisation linéaire ultérieure et soit \(S\) la permutation
\((Sx)_i=x_{i+1}\), avec des indices modulo \(12\).
Les lecteurs \(I-S^3\) et \(I-S^4\) sont ceux de l’extracteur transversal
de \cref{img09:imagen}.
La compression nonadique de \cref{nuclear:10} fournit,
pour \(j=3,4\),



\begin{equation}
T_j=\frac{8I+S^j}{9},\qquad D_j=\frac{\sqrt8}{9}(S^j-I),
\qquad T_j^*T_j+D_j^*D_j=I.
\label{nuc11:eq:1}
\end{equation}



La partition \(8+1\) et la normalisation \(9\) proviennent des neuf
positions et de leur contraste collectif.
On considère la composition explicite de lecteurs suivante,
appliquée au registre engendré :



\begin{equation}
m_0=D_3K,\qquad m_1=D_4T_3K,
\qquad M K=(m_0,m_1).
\label{nuc11:eq:2}
\end{equation}



La seconde mémoire agit sur la lecture produite par la première
étape. Les compléments sont archivés séparément.
L’évolution complète avec réentrée correspond au couplage
réversible de \cref{nuclear:10}.

\subsubsection{Inversa exacta}

Puisque \((S^3)^4=I\),



\begin{equation}
T_3^{-1}=\frac{512I-64S^3+8S^6-S^9}{455}.
\label{nuc11:eq:3}
\end{equation}



En effet,
\((8I+S^3)(512I-64S^3+8S^6-S^9)=4095I\)
et \(4095=9\cdot455\).
Par conséquent, les mémoires fournissent



\begin{equation}
b=-\frac9{\sqrt8}m_0=(I-S^3)K,
\quad c=-\frac9{\sqrt8}T_3^{-1}m_1=(I-S^4)K.
\label{nuc11:eq:4}
\end{equation}



La commutation de \(T_3\) avec \(S^4\) justifie la seconde égalité.
Définissons



\[
w_i=c_i-b_{i+1},\quad B_0=0,\quad
B_i=\sum_{j<i}w_j.
\]



L’identité \(w_i=K_i-K_{i+1}\) prouve que \(B_i=K_0-K_i\).
Si \(q=\sum_{i=0}^{11}K_i\), on récupère exactement




\begin{equation}
K_i=\frac{q+\sum_{j=0}^{11}B_j}{12}-B_i.
\label{nuc11:eq:5}
\end{equation}



C’est l’inverse \eqref{img09:inversa}, composé maintenant avec la
compression et la mémoire. Avec \(q=0\), la formule renvoie
\(P_{11}K\), où



\[
P_{11}=I-\frac1{12}\mathbf1\mathbf1^*.
\]



Les deux mémoires retrouvent les onze composantes de moyenne nulle.
La charge uniforme se retrouve au moyen de \(q\).
Le noyau de \(M\) est précisément la droite uniforme :
annuler les deux différences implique la fixité par \(S^3\) et \(S^4\) ;
comme \(S=S^4(S^3)^{-1}\), cela implique la fixité par \(S\).

\subsubsection{Borne optimale de stabilité
}

Le gramien du procédé est




\begin{equation}
G=M^*M=D_3^*D_3+T_3^*D_4^*D_4T_3
=I-(T_4T_3)^*(T_4T_3).
\label{nuc11:eq:6}
\end{equation}



Dans le mode \(k\) de \(S\), la valeur propre de \(T_j^*T_j\) est

\[
 \frac{65+16\cos(2\pi jk/12)}{81}.
\]
Sa substitution dans la matrice de Gram précédente donne, avec multiplicités,



\begin{equation}
\operatorname{Spec}G=\left\{
0^{[1]},\ (16/81)^{[2]},\ (8/27)^{[2]},\ (32/81)^{[1]},
(952/2187)^{[4]},\ (1256/2187)^{[2]}\right\}.
\label{nuc11:eq:7}
\end{equation}



Le mode nul est la moyenne. Soit \(L\) l’inverse aux moindres carrés
de \(M\) sur \(\mathbf1^\perp\), étendu à l’espace des
registres par




\[
L=(G|_{\mathbf1^\perp})^{-1}M^*.
\]



Alors



\begin{equation}
LM=P_{11},\qquad \|L\|=\frac94.
\label{nuc11:eq:8}
\end{equation}



La borne est optimale parce que les modes \(k=3,9\) atteignent \(16/81\).
Pour des données perturbées et une reconstruction avec \(L\),



\begin{equation}
\|\delta K\|^2\le \frac{|\delta q|^2}{12}
+\frac{81}{16}\bigl(\|\delta m_0\|^2+\|\delta m_1\|^2\bigr).
\label{nuc11:eq:9}
\end{equation}



Les composantes uniforme et centrée sont orthogonales. La formule \eqref{nuc11:eq:5} constitue un inverse exact sur des données compatibles ; \eqref{nuc11:eq:8} spécifie l’extension stable utilisée pour des données bruitées, qui ne satisfont pas nécessairement les compatibilités de \eqref{nuc11:eq:5}.

\subsubsection{Exemple avec le registre exprimé}

Le registre de \eqref{eq:k-vector} est



\begin{equation}
K=(234,543,140,729,659,824,621,58,914,794,146,601).
\label{nuc11:eq:10}
\end{equation}



C’est une sortie du générateur précédent, utilisée comme cas
d’application. Les coefficients et les lecteurs sont fixés avant
d’évaluer ce registre. Pour stocker les mémoires en arithmétique
rationnelle, nous écrivons \(z_i=m_i/\sqrt8\).
Les deux registres effectifs sont




\[
9z_0=(495,116,684,-108,-601,90,173,88,-313,-560,397,-461),
\]





\begin{equation}
81z_1=(2729,2503,3818,-5843,2583,-920,-4051,4338,-5312,-1583,233,1505).
\label{nuc11:eq:11}
\end{equation}



La substitution de ces registres dans les formules de reconstruction,
avec \(q=6263\), reproduit les douze composantes de \(K\).
Avec \(q=0\), elle reproduit \(K-(6263/12)\mathbf1\).
Le certificat supplémentaire reconstruit d’abord le registre à partir
des mémoires, puis le compare à \eqref{eq:k-vector}.

Ajouter une troisième mémoire \(D_3T_4T_3K\) élève la plus petite valeur
propre positive du Gram à \(8/27\).
L’amplification optimale descend de \(9/4\) à \(\sqrt{27/8}\).
Une troisième mémoire \(D_4T_4T_3K\) maintient \(16/81\).
Le choix entre la répétition du pas \(3\) ou du pas \(4\) a
un effet de stabilité calculable avant l’évaluation du registre.

\subsection{Reconstruction du projecteur dimensionnel à partir de deux mémoires
}

\subsubsection{Composition avec l’incidence déjà construite}

Le système de coordonnées exceptionnel de \cref{exc:k-direccion} définit la
représentation orthogonale de \(A_5\) sur les douze classes de
\(A_5/C_5\), avec des représentants et un caractère explicites.
Son projecteur isotypique tridimensionnel est
\eqref{eq:k-proyectores-incidencia} :



\begin{equation}
P_3=\frac3{60}\sum_{a\in A_5}\chi_3(a^{-1})\rho(a),
\quad P_3P_{11}=P_3,\quad P_3\mathbf1=0.
\label{nuc11:eq:12}
\end{equation}



Le registre engendré précède cette représentation.
Le \cref{lem:k-sector-no-nulo} et
\eqref{eq:k-norma-direccion} et \eqref{eq:k-proyector-diez} établissent



\[
u_K=P_3P_{11}K,\qquad
\|u_K\|^2=\frac{6638585+2275584\sqrt5}{20}>0,
\]





\begin{equation}
P_{10}(K)=P_{11}-\frac{u_Ku_K^*}{\|u_K\|^2}.
\label{nuc11:eq:13}
\end{equation}



La nouvelle composition de \eqref{nuc11:eq:8}, \eqref{nuc11:eq:12} et \eqref{nuc11:eq:13} est



\begin{equation}
\boxed{
u_K=P_3L(m_0,m_1),\qquad
P_{10}(K)=P_{11}-
\frac{P_3L(m)\,[P_3L(m)]^*}{\|P_3L(m)\|^2}.
}
\label{nuc11:eq:14}
\end{equation}



\begin{proof}
Les identités \(LM=P_{11}\) et \(P_3P_{11}=P_3\) donnent
\(P_3LMK=P_3K=u_K\).
La non-nullité est évaluée dans \eqref{eq:k-norma-direccion}.
Le quotient de rang un coïncide donc avec celui de
\eqref{eq:k-proyector-diez}.
La composition récupère le projecteur orthogonal intégral de rang
dix à partir des deux mémoires, indépendamment de la charge
\(q\) et de la lecture terminale \(T_4T_3K\).
\end{proof}

La réduction \(12\to11\) élimine la direction uniforme et la
réduction \(11\to10\) utilise le vecteur sélectionné par \(K\).
La composition récupère la matrice complète du projecteur dans le
support linéaire d’incidence.
Ses réalisations dimensionnelles maintiennent les applications de
\cref{prop:k-reduccion-diez}.


\subsubsection{Erreur du projecteur}

Soient \(\varepsilon\) la norme conjointe de l’erreur des deux
mémoires et \(\widehat u=P_3L(m+\delta m)\).
Alors \(\|\widehat u-u_K\|\le9\varepsilon/4\).
Si \(9\varepsilon/4<\|u_K\|\), le vecteur \(\widehat u\) est non
nul et



\begin{equation}
\boxed{
\|\widehat P_{10}-P_{10}\|
\le\frac{9\varepsilon}{4\|u_K\|}.
}
\label{nuc11:eq:15}
\end{equation}



\begin{proof}
La distance en norme d’opérateur entre deux projecteurs orthogonaux
sur des droites est le sinus de l’angle entre elles.
La distance de \(u_K\) à la droite de \(\widehat u\) est au plus
\(\|u_K-\widehat u\|\). Diviser par \(\|u_K\|\) prouve la borne.
Ici, \(\|u_K\|\approx765{,}733\), et le coefficient est
environ \(0{,}00293836\).
\end{proof}
Pour des registres proches de \(\ker P_3\), la stabilité de la
direction dépend de \(\|u_K\|\).
La borne \(9/4\) pour le registre centré reste valable.


\subsubsection{Un descripteur scalaire de K, avec invariances explicites}

On distingue le registre \(K\), ses lectures scalaires et ses
invariants. Pour \(K\notin\RR\mathbf1\), nous définissons




\begin{equation}
\eta_K=\frac{\|P_3P_{11}K\|^2}{\|P_{11}K\|^2}.
\label{nuc11:eq:16}
\end{equation}



Pour le registre \eqref{eq:k-vector},
\(\sum_iK_i^2=4251257\) et
\(\|P_{11}K\|^2=11789915/12\). Ainsi,



\begin{equation}
\eta_K=
\frac{3(6638585+2275584\sqrt5)}{5\cdot11789915}
=\frac{13089003+13653504\varphi_{\rm HMT}}{58949575}
\approx0,5967954228259091.
\label{nuc11:eq:17}
\end{equation}



La seconde écriture utilise la reconnaissance quadratique ultérieure
de \(\varphi_{\rm HMT}\).
On a \(0\le\eta_K\le1\).
Le quotient est invariant sous
\(K\mapsto aK+b\mathbf1\), avec \(a\ne0\), et sous des changements
orthogonaux simultanés de \(K,P_{11},P_3\).
Dans le système de coordonnées fixe, il est invariant sous \(\rho(A_5)\), puisque les deux
projecteurs commutent avec cette représentation.
Il est récupérable à partir des mémoires parce que les deux normes qui le
définissent le sont.
Il mesure la fraction quadratique du registre centré située dans le
secteur isotypique choisi.
Le registre complet conserve en outre l’information que ce
descripteur scalaire identifie.
La provenance de cette formalisation supplémentaire est documentée dans
le supplément ; les invariances affines prouvées ici appartiennent
à la réalisation linéaire.

\subsection{Application sur un vecteur effectif de la construction exceptionnelle
}

La construction de
\cref{km:orden-tres-reticular,km:estabilidad-vecino}
utilise douze fibres \(A_2\), leur recollement et le voisin réticulaire.
Le mot tritique est
\(c^\star=(1,1,1,1,1,1,2,1,1,1,1,1)\).
Ses blocs agissent par \(C^{-1}\) lorsque la composante est \(1\)
et par \(C\) lorsqu’elle est \(2\), dans les systèmes de coordonnées spécifiés par
\cref{km:palabra-total,km:isometria-ternaria}.
L’opérateur \(g\) résultant satisfait
\(g^3=I\), \(g^2+g+I=0\), et est orthogonal pour la somme des
métriques \(A_2\).
La descente au recollement et au voisin est démontrée dans
\cref{km:estabilidad-vecino}.
Les calculs suivants agissent sur son support réel de
dimension \(24\).

Dans chaque fibre,

\[
 C=\begin{pmatrix}0&-1\\1&-1\end{pmatrix},
 \qquad
 G_{A_2}=\begin{pmatrix}2&-1\\-1&2\end{pmatrix}.
\]
Sur la somme des douze fibres, on définit




\[
T=\frac{8I+g}{9},\qquad D=\frac{\sqrt8}{9}(g-I).
\]



Avec l’adjoint relatif à cette métrique,
\(g^*=g^{-1}=-g-I\), et on obtient




\begin{equation}
T^*T=\frac{19}{27}I,
\qquad D^*D=\frac8{27}I,
\qquad D^{-1}=-\frac3{\sqrt8}(g+2I).
\label{nuc11:eq:18}
\end{equation}



L’inverse provient de \((g-I)(g+2I)=-3I\).
La construction du voisin fixe le vecteur de
\eqref{exc:vector},
\(v=4\rho_1+\rho_2+\cdots+\rho_{12}\), avec une racine
\(\rho_b\) de norme quadratique \(2\) dans chaque fibre.
Par conséquent,




\begin{equation}
\|v\|^2=54,\qquad \|Tv\|^2=38,\qquad \|Dv\|^2=16,
\qquad v=-\frac3{\sqrt8}(g+2I)Dv.
\label{nuc11:eq:19}
\end{equation}



Un seul complément récupère ce vecteur complet.
La lecture conserve \(38\) unités de norme quadratique et le
registre contient \(16\) ; leur somme est \(54\).
Le résultat utilise le vecteur de \eqref{exc:vector} et ses fibres
réticulaires.


\subsection{Loi de conservation à chaque degré extérieur}

\subsubsection{Énoncé et preuve}

Soient \(P_0=I\), \(P_n=T_{n-1}\cdots T_0\) et




\[
\mathcal J_Nv=(P_Nv,D_0P_0v,\ldots,D_{N-1}P_{N-1}v).
\]



La loi de \cref{nuclear:8,nuclear:10} conserve la forme \(G_0\)
au moyen de la somme orthogonale des formes correspondantes :
\(\mathcal J_N^*G_{\rm sal}\mathcal J_N=G_0\).
Pour chaque degré extérieur fini admissible \(p\), on en déduit



\begin{equation}
\boxed{
(\Lambda^p\mathcal J_N)^*(\Lambda^pG_{\rm sal})
(\Lambda^p\mathcal J_N)=\Lambda^pG_0.
}
\label{nuc11:eq:20}
\end{equation}



La forme induite est définie par



\[
\langle x_1\wedge\cdots\wedge x_p,
y_1\wedge\cdots\wedge y_p\rangle_{\Lambda^pG}
=\det(\langle x_i,y_j\rangle_G)_{i,j=1}^p.
\]



\begin{proof}
Substituer le bilan de degré un dans chaque entrée de ce déterminant. L’égalité est satisfaite sur les vecteurs décomposables et s’étend bilinéairement parce que ceux-ci engendrent la puissance extérieure.

\end{proof}

La décomposition pertinente est



\[
\Lambda^p\Bigl(\bigoplus_aH_a\Bigr)
\simeq\bigoplus_{\sum n_a=p}\bigotimes_a\Lambda^{n_a}H_a.
\]



Dans le cas positif, si
\(\mathcal J_{p,\mathbf n}\) désigne les composantes de cette application,



\begin{equation}
\boxed{\|\xi\|^2=\sum_{\sum n_a=p}
\|\mathcal J_{p,\mathbf n}\xi\|^2.}
\label{nuc11:eq:21}
\end{equation}



Les composantes purement visibles, purement
enregistrées et mixtes sont conservées.
Pour \(p=2\) et \(p=3\), les normes
quadratiques d’aire et de volume sont respectivement conservées.
Pour une forme alternée ou indéfinie, l’identité extérieure
conserve son interprétation bilinéaire.

\subsubsection{Distribution exacte entre les secteurs extérieurs}

Dans la réalisation de Coxeter précédente,
\(a=19/27\) et \(b=8/27\).
Normaliser \(T\) et \(D\) donne deux isométries, de sorte que la
composante ayant \(k\) directions dans la mémoire a pour norme quadratique




\begin{equation}
\boxed{
\|\mathcal J_{p,k}\xi\|^2=
\binom pk a^{p-k}b^k\|\xi\|^2.
}
\label{nuc11:eq:22}
\end{equation}



Cela se prouve en développant extérieurement
\(x\mapsto(\sqrt a\,x,\sqrt b\,x)\).
Les composantes ayant des positions de mémoire distinctes sont
orthogonales ; les isométries normalisées \(T\) et \(D\)
conservent leurs normes.

Pour un bivecteur de norme un,

\begin{center}
\begin{tabular}{@{}lr@{}}
\toprule
Sector & Norme quadratique
\\
\midrule
Deux directions dans la lecture & \(361/729\)\\
Une direction dans la lecture et une autre dans la mémoire & \(304/729\)\\
Deux directions en mémoire
 & \(64/729\)\\
\bottomrule
\end{tabular}
\end{center}

La somme vaut un. Le secteur mixte contient exactement \(304/729\),
soit environ \(41{,}7\%\), de l’identité d’aire au carré.
En degré trois, les quatre composantes sont
\((6859,8664,3648,512)/19683\) ; les composantes mixtes ont pour somme
\(12312/19683\).


La loi s’applique à \(0\le p\le24\) dans le porteur des douze
fibres \(A_2\), et fonctoriellement à d’autres porteurs engendrés.
Le degré extérieur, la résolution de raffinement, l’échelle d’une
métrique et la dimension physique ont des typages distincts.
Le changement \(G\mapsto s^2G\), avec \(s>0\), multiplie la norme
quadratique de degré \(p\) par \(s^{2p}\) ; les deux membres de
l’identité extérieure se transforment de la même manière.
La réalisation des dimensions \(10\) et \(11\) conserve les
applications dimensionnelles préalablement spécifiées.

\subsubsection{Naturalité et profondeur arbitraire}

Le raffinement des histoires engendré conserve la mesure par



\[
R\delta_h=\sum_{\epsilon\in\operatorname{Out}(h)}
\sqrt{p_h(\epsilon)}\,\delta_{h\epsilon}.
\]



Les probabilités proviennent des multiplicités et de la
normalisation des émissions APP--TRIT--TPK de
\cref{iv:cont:medida,nuclear:1}.
Le registre identifie l’état précédent de chaque émission.
Si le carré de transport est
\(\mathcal J'_NR=(\bigoplus_aR_a)\mathcal J_N\), alors



\begin{equation}
(\Lambda^pJ'_N)(\Lambda^pR)
=\Lambda^p(\bigoplus R_a)(\Lambda^pJ_N).
\label{nuc11:eq:23}
\end{equation}



On applique
\(\Lambda^p(AB)=\Lambda^p(A)\Lambda^p(B)\).
Les carrés des cinq opérations corrélatives et leur limite
commune sont ceux de
\cref{iv:rectangular-natural,iv:cont:memoria-natural,iv:cont:mapa-correlativo} ; toutes agissent dans la même construction
du continu.


Dans le cas positif non stationnaire,
\cref{nuclear:10} démontre
\(P_N^*P_N\downarrow A_\infty\) fortement et



\[
\mathcal Jv=(A_\infty^{1/2}v,(D_nP_nv)_{n\ge0}),
\qquad \mathcal J^*\mathcal J=I.
\]



Par densité, \(\Lambda^p\mathcal J\) est isométrique pour chaque
degré fini \(p\).
La somme extérieure s’étend à des multi-indices de support fini
sur la somme de Hilbert dénombrable ; Parseval justifie le passage
à la limite.
Le terminal \(A_\infty^{1/2}\) et les mémoires conservent
conjointement la forme initiale à profondeur arbitraire.


\subsubsection{Précision de reconstruction des aires et des volumes}

Pour le lecteur \(M\), ordonnons les onze valeurs propres
positives de \(G\), avec leur multiplicité, sous la forme
\(\lambda_1\le\cdots\le\lambda_{11}\).
Aux degrés \(1\le p\le11\) de \(\mathbf1^\perp\),
la décomposition singulière induit



\begin{equation}
\|(\Lambda^pM)^\dagger\|
=(\lambda_1\cdots\lambda_p)^{-1/2}.
\label{nuc11:eq:24}
\end{equation}



Les modes extérieurs utilisent \(p\) directions distinctes.
Le spectre calculé donne les constantes optimales
\(9/4\), \(81/16\), \(243\sqrt6/64\) et \(2187/128\)
pour \(p=1,2,3,4\).
En degré douze, \(\Lambda^{12}M=0\) ; la récupération du volume
complet incorpore aussi la charge uniforme.
Le produit tensoriel de \(p\) lecteurs indépendants sur
la moyenne nulle a pour borne \((9/4)^p\).
Tensoriser le lecteur par une identité conserve la borne \(9/4\).
Le registre complet \(\mathcal J\) a un inverse de norme un
sur son image à chaque degré.

\subsection{Mémoire du transport et holonomie relative}

Le couplage réversible de \cref{nuclear:10} est



\[
U(B,C)=\frac19
\begin{pmatrix}8B+C&\sqrt8(C-B)\\
\sqrt8(C-B)&B+8C\end{pmatrix}
=Q^*\operatorname{diag}(B,C)Q,
\]



avec \(Q\) constante et orthogonale, fixée par la partition \(8+1\).
Par conséquent,
\(U(B_2,C_2)U(B_1,C_1)=U(B_2B_1,C_2C_1)\),
en conservant l’ordre des facteurs.
Sur une boucle \(\gamma\), soient \(H_B,H_C\) les produits
ordonnés des deux transports.
Le bloc de mémoire de l’évolution complète est




\[
D_\gamma=\frac{\sqrt8}{9}(H_C-H_B).
\]



On reconstruit exactement



\begin{equation}
\boxed{H_B^{-1}H_C=I+\frac9{\sqrt8}H_B^{-1}D_\gamma.}
\label{nuc11:eq:25}
\end{equation}



L’égalité est opératorielle : enregistrer la réponse du bloc
sur une base retrouve l’holonomie relative complète.
Son évaluation sur un vecteur retrouve l’action correspondante.
Sous un changement de système de coordonnées au point de base, les deux membres sont conjugués.
Si \(H_B=I\), le complément détermine \(H_C-I\).
Ainsi sont reliés la mémoire et le défaut de transport sur un
circuit ; la réalisation différentielle utilise en outre
sa connexion et la limite des circuits.

\subsubsection{Une boucle elliptique–parabolique produit une auto-échelle hyperbolique
}

La construction de \cref{euler-trit-generadores} contient le
Coxeter \(C\) et le nilpotent

\[
 N_0=\begin{pmatrix}0&1\\0&0\end{pmatrix}
\]
dans leurs porteurs respectifs.
On considère leurs réalisations dans un plan orienté commun,
au moyen des bases indiquées et de la forme
\[
 \Omega=\begin{pmatrix}0&1\\-1&0\end{pmatrix}.
\]
Le Coxeter \(C\) et le transport \(I+N_0\) ont un déterminant
un et conservent \(\Omega\).
Ce système de coordonnées permet de les composer en conservant la distinction entre
les générateurs de leurs régimes et leurs métriques positives propres.

Soit
\[
 P=I+N_0=\begin{pmatrix}1&1\\0&1\end{pmatrix}.
\]
Comme \(N_0^2=0\), on a \(P^n=I+nN_0\) pour tout
\(n\in\ZZ\).
Le lacet ordonné de quatre transports donne



\begin{equation}
H_n=C^{-1}P^{-n}CP^n
=\begin{pmatrix}1+n&n^2\\n&n^2-n+1\end{pmatrix}.
\label{nuc11:eq:28}
\end{equation}



Pour \(n\ne0\), sa différence normalisée est



\[
F_n=\frac{H_n-I}{n}
=\begin{pmatrix}1&n\\1&n-1\end{pmatrix}.
\]



En multipliant ces matrices, on obtient, sans introduire de valeurs spectrales,



\begin{equation}
\boxed{F_n^2=H_n,\qquad F_n^2-nF_n-I=0.}
\label{nuc11:eq:29}
\end{equation}



Avec \(H_B=I\), la mémoire est
\(D_{\gamma,n}=(\sqrt8/9)nF_n\).
Par conséquent,



\begin{equation}
\boxed{H_n=\left(\frac9{n\sqrt8}D_{\gamma,n}\right)^2.}
\label{nuc11:eq:30}
\end{equation}



La réponse matricielle de mémoire reconstruit l’auto-échelle
hyperbolique. Pour \(n>0\), la reconnaissance spectrale ultérieure
identifie



\begin{equation}
\rho_n=\frac{n+\sqrt{n^2+4}}2,\quad
\operatorname{Spec}F_n=\{\rho_n,-\rho_n^{-1}\},\quad
\operatorname{Spec}H_n=\{\rho_n^2,\rho_n^{-2}\}.
\label{nuc11:eq:31}
\end{equation}



Le cas \(n=1\) produit le polynôme d’or.
La relation avec l’opérateur
\(F_{\rm av}=
\begin{pmatrix}0&1\\1&1\end{pmatrix}\)
engendré par le calendrier dans \eqref{eq:fav} conserve le changement
de système de coordonnées :



\begin{equation}
F_1=PF_{\rm av}P^{-1},\qquad
H_1=PF_{\rm av}^{\,2}P^{-1}.
\label{nuc11:eq:32}
\end{equation}



Le transport \(P\), de déterminant un, entrelace les deux
matrices et conserve \(\Omega\).
Le cas \(n=2\) produit les lectures spectrales
\(1+\sqrt2\) et \(3+2\sqrt2\).
La composition relie ces réalisations en conservant leurs
généalogies respectives ; la sélection d’un mot de quatre
facteurs reste spécifiée comme partie du lecteur.


Dans le même exemple \(n=1\), on obtient le bilan orienté



\begin{equation}
D_\gamma^{\mathsf T}\Omega D_\gamma=-\frac8{81}\Omega,
\qquad
T_\gamma^{\mathsf T}\Omega T_\gamma=\frac{89}{81}\Omega.
\label{nuc11:eq:33}
\end{equation}



La somme des deux contributions est \(\Omega\).
Le terme de mémoire a l’orientation opposée :
la contribution visible \(89/81\) est compensée par \(-8/81\).
Ce bilan alterné et la répartition positive
\(19/27+8/27=1\) conservent leurs formes respectives.

L’expression de \(H_n\) vaut pour tout \(n\in\ZZ\).
La définition de \(F_n\) et la reconstruction au moyen de son carré
utilisent \(n\ne0\) ; la paramétrisation spectrale précédente utilise
\(n>0\).
En général,
\[
 D_{\gamma,n}^{\mathsf T}\Omega D_{\gamma,n}
 =-\frac{8n^2}{81}\Omega,\qquad
 T_{\gamma,n}^{\mathsf T}\Omega T_{\gamma,n}
 =\left(1+\frac{8n^2}{81}\right)\Omega.
\]
Les identités \(N_0^2=0\) et \(C^2+C+I=0\) démontrent la
collection complète ; le programme supplémentaire vérifie ses
spécialisations exactes.
La provenance de cette formalisation supplémentaire est documentée
dans le supplément.

\subsection{Symétrie, action et charge conservée}

Soient \((E_n)_{n\in\ZZ}\) des espaces vectoriels réels de dimension
finie et \(U_n:E_n\to E_{n+1}\) les transports inversibles de
l’évolution complète.
Le relèvement cotangent agit sur
\(q_n\in E_n\) et \(p_n\in E_n^*\).
Dans des bases duales, nous définissons le lagrangien discret auxiliaire

\[
 L_n(q_n,q_{n+1},p_{n+1})
 =p_{n+1}^{\mathsf T}(q_{n+1}-U_nq_n).
\]
L’action est utilisée au sens local :
pour chaque variation à support fini, on somme \(L_n\) sur un
intervalle fini contenant tous les termes affectés.
La première variation est indépendante de tout agrandissement
de cet intervalle. Concrètement,
\begin{equation}
 \delta\mathscr S
 =\sum_{n\in\ZZ}
 \left[
  \delta p_n^{\mathsf T}(q_n-U_{n-1}q_{n-1})
  +\delta q_n^{\mathsf T}(p_n-U_n^{\mathsf T}p_{n+1})
 \right],
 \label{nuc11:accion-local}
\end{equation}
où la somme est finie pour chaque variation admise.
La stationnarité par rapport à toutes ces variations équivaut à
\[
 q_{n+1}=U_nq_n,\qquad
 p_n=U_n^{\mathsf T}p_{n+1}
 \quad(n\in\ZZ).
\]
La seconde équation est
\(p_{n+1}=U_n^{-\mathsf T}p_n\), c’est-à-dire le relèvement
cotangent du transport.


\begin{proposition}[Charge variationnelle du transport entrelacé]
\label{nuc11:noether}
Soit \(A_n\in\operatorname{End}(E_n)\) une collection qui satisfait
\[
 A_{n+1}U_n=U_nA_n.
\]
Sur les solutions stationnaires de
\eqref{nuc11:accion-local}, la charge
\begin{equation}
 \boxed{J_n=p_n^{\mathsf T}A_nq_n=J_{n+1}}
 \label{nuc11:carga}
\end{equation}
est constante.
\end{proposition}
\begin{proof}
Les équations d’évolution et l’entrelacement donnent

\[
 J_{n+1}
 =p_{n+1}^{\mathsf T}A_{n+1}U_nq_n
 =p_{n+1}^{\mathsf T}U_nA_nq_n
 =p_n^{\mathsf T}A_nq_n.
\]
La conservation possède en outre une dérivation variationnelle explicite.
Pour un paramètre constant \(\varepsilon\), les transformations
\[
 q_n\mapsto e^{\varepsilon A_n}q_n,\qquad
 p_n\mapsto e^{-\varepsilon A_n^{\mathsf T}}p_n
\]
conservent chaque \(L_n\), car
\(e^{\varepsilon A_{n+1}}U_n=U_ne^{\varepsilon A_n}\).
Si le paramètre est remplacé par une suite
\((\varepsilon_n)\) à support fini, la première variation est
\[
 \delta L_n
 =(\varepsilon_{n+1}-\varepsilon_n)
   p_{n+1}^{\mathsf T}U_nA_nq_n.
\]
Sur une solution, le dernier facteur est \(J_n\).
La sommation finie par parties donne

\[
 \delta\mathscr S
 =\sum_{n\in\ZZ}\varepsilon_n(J_{n-1}-J_n).
\]
La stationnarité pour toute suite à support fini produit
\eqref{nuc11:carga}.
\end{proof}

Pour les lecteurs cycliques précédents, on prend

\[
 A_0=S-S^{-1},\qquad
 \mathcal A=\operatorname{diag}(A_0,A_0)
\]
dans le porteur complet \(V\oplus V\).
La matrice \(\mathcal A\) commute avec \(U(I,S^3)\) et
\(U(I,S^4)\), car ses blocs sont des polynômes en \(S\).
Par conséquent, \(p_n^{\mathsf T}\mathcal A q_n\) conserve
une corrélation orientée entre positions voisines.

La construction spécifie un lagrangien, sa symétrie et la charge
conservée correspondante, au sens variationnel développé
dans \cref{nuclear:7}.
Son domaine est la réalisation cotangente des lecteurs.
L’action physique et ses unités conservent les opérateurs et les
normalisations des constructions correspondantes.
Le lagrangien auxiliaire s’obtient après le transport \(U_n\).
